\chapter{Low-Power, Touch and Crypto}
\label{ch:hal3}

These drivers need no finalization, so they build under \emph{every} runtime
profile.

\section{RTC: deep sleep and retained memory}\index{RTC driver}

The RTC domain keeps a small slow-memory block alive across deep sleep, and
drives the wake timer. \texttt{Read}/\texttt{Write} access retained 32-bit words;
\texttt{Deep\_Sleep\_For}/\texttt{Until} enter deep sleep (the chip resets on
wake, so execution resumes at the top of \texttt{Main}); \texttt{Last\_Wake}
reports why it woke.

\begin{lstlisting}
function  Read  (Index : Word_Index) return Interfaces.Unsigned_32;
procedure Write (Index : Word_Index; Value : Interfaces.Unsigned_32);
function  Last_Wake return Wake_Cause;   --  Power_On / Deep_Sleep_Timer / ...
procedure Deep_Sleep_For   (Wake_After : Duration);
procedure Deep_Sleep_Until (Pin : Pin_Id; High : Boolean := True);
\end{lstlisting}

\textbf{Parameters.}
\begin{description}[style=nextline]
\item[\texttt{Index}] retained-word slot, \texttt{0\,..\,2047}
  (\texttt{Word\_Index}, the 8\,KB RTC slow memory as 32-bit words).
\item[\texttt{Value}] the 32-bit word to store (\texttt{Unsigned\_32}).
\item[\texttt{Wake\_After}] (\texttt{Deep\_Sleep\_For}) the sleep length as an Ada
  \texttt{Duration} in \emph{seconds}; approximate, since the RTC slow clock is
  uncalibrated.
\item[\texttt{Pin}, \texttt{High}] (\texttt{Deep\_Sleep\_Until}) the wake pad and
  the level to wake on; \texttt{High => True} (default) wakes on a high level.
\item[\texttt{Last\_Wake}] returns \texttt{Power\_On}, \texttt{Deep\_Sleep\_Timer},
  \texttt{Deep\_Sleep\_GPIO} or \texttt{Other\_Reset}.
\end{description}

\lstcap{Example 1: a boot counter that survives deep sleep.}
\begin{lstlisting}
with Interfaces;     use Interfaces;
with Ada.Text_IO;    use Ada.Text_IO;
with ESP32S3.RTC;    use ESP32S3.RTC;
procedure RTC_Boot is
   N : Interfaces.Unsigned_32;
begin
   if Last_Wake = Deep_Sleep_Timer then
      N := Read (0) + 1;     --  woke from our own timer
   else
      N := 1;                --  cold boot
   end if;
   Write (0, N);
   Put_Line ("boot #" & N'Image);
   Deep_Sleep_For (5.0);     --  Duration in seconds; wakes and re-enters Main
end RTC_Boot;
\end{lstlisting}

\lstcap{Example 2: wake on a button.}
\begin{lstlisting}
with ESP32S3.RTC; use ESP32S3.RTC;
procedure RTC_Wake is
begin
   Deep_Sleep_Until (Pin => 0, High => False);  --  sleep until GPIO0 goes low
end RTC_Wake;
\end{lstlisting}
(After a deep-sleep wake the driver also disarms the super-watchdog
(\texttt{Disable\_Super\_Watchdog}) so a staying-awake app is not reset.)

\section{RTC\_IO: pad hold and RTC pulls}\index{RTC\_IO driver}

Latch a pad's level so it survives \texttt{GPIO} writes and deep sleep, and apply
RTC-domain pulls.

\begin{lstlisting}
procedure Hold (Pin : RTC_Pin);   procedure Release (Pin : RTC_Pin);
function  Is_Held (Pin : RTC_Pin) return Boolean;
procedure Set_Pull (Pin : RTC_Pin; Mode : Pull_Mode);   --  No_Pull / Up / Down
\end{lstlisting}

\textbf{Parameters.}
\begin{description}[style=nextline]
\item[\texttt{Pin}] an RTC-capable pad (\texttt{RTC\_Pin}), restricted to
  \texttt{0\,..\,21} (only those reach the RTC domain).
\item[\texttt{Mode}] (\texttt{Set\_Pull}) the RTC-domain pull: \texttt{No\_Pull},
  \texttt{Up} or \texttt{Down}.
\end{description}

\lstcap{Example 1: hold a relay line high through sleep.}
\begin{lstlisting}
with ESP32S3.GPIO;
with ESP32S3.RTC_IO; use ESP32S3.RTC_IO;
procedure RTCIO_Hold is
begin
   ESP32S3.GPIO.Configure (5, Mode => ESP32S3.GPIO.Output);
   ESP32S3.GPIO.Set (5);
   Hold (5);                 --  GPIO writes now ignored; level frozen
   --  ... deep sleep ...    --  pad stays high across sleep
   Release (5);              --  GPIO control restored
end RTCIO_Hold;
\end{lstlisting}

\lstcap{Example 2: an RTC pull-up on a wake pin.}
\begin{lstlisting}
with ESP32S3.RTC_IO; use ESP32S3.RTC_IO;
procedure RTCIO_Pull is
begin
   Set_Pull (0, Up);
end RTCIO_Pull;
\end{lstlisting}

\section{Touch: capacitive sensing}\index{Touch driver}

\begin{lstlisting}
procedure Setup;                       --  bring up the touch FSM (call once)
procedure Enable (Ch : Channel);       --  channel n is on GPIO n (n = 1 .. 14)
function  Read (Ch : Channel) return Natural;   --  raw self-capacitance count
function  Touched (Ch : Channel; Reference : Natural; Margin : Natural := 20_000)
   return Boolean;
\end{lstlisting}

\textbf{Parameters.}
\begin{description}[style=nextline]
\item[\texttt{Ch}] touch channel \texttt{1\,..\,14} (\texttt{Channel}); channel
  \emph{n} is wired to GPIO \emph{n}.
\item[\texttt{Read}] returns a raw self-capacitance count (\texttt{Natural};
  larger means more capacitance; \texttt{0} means not measured yet).
\item[\texttt{Reference}] (\texttt{Touched}) the untouched baseline count you
  captured earlier with \texttt{Read}.
\item[\texttt{Margin}] how far the live count must deviate from \texttt{Reference}
  to register a touch; default \texttt{20\_000}.
\end{description}

\lstcap{Example 1: read two pads' baselines.}
\begin{lstlisting}
with Ada.Text_IO;   use Ada.Text_IO;
with ESP32S3.Touch; use ESP32S3.Touch;
procedure Touch_Read_Base is
begin
   Setup;  Enable (1);  Enable (3);
   Put_Line ("ch1 =" & Natural'Image (Read (1)) &
             "  ch3 =" & Natural'Image (Read (3)));
end Touch_Read_Base;
\end{lstlisting}

\lstcap{Example 2: detect a finger against a captured baseline.}
\begin{lstlisting}
with Ada.Text_IO;   use Ada.Text_IO;
with ESP32S3.Touch; use ESP32S3.Touch;
procedure Touch_Detect is
   Base : constant Natural := Read (1);   --  with nothing touching
begin
   loop
      if Touched (1, Base, Margin => 50_000) then
         Put_Line ("touched!");
      end if;
   end loop;
end Touch_Detect;
\end{lstlisting}

\section{SHA: hardware hashing}\index{SHA driver}

\begin{lstlisting}
function Hash_1   (Data : Byte_Array) return SHA1_Digest;    --  20 bytes
function Hash_224 (Data : Byte_Array) return SHA224_Digest;  --  28 bytes
function Hash_256 (Data : Byte_Array) return SHA256_Digest;  --  32 bytes
\end{lstlisting}

\textbf{Parameters.} Each \texttt{Hash\_*} takes one \texttt{Data : Byte\_Array} of
\emph{any} length and returns a fixed-size digest: \texttt{SHA1\_Digest}
(20 bytes), \texttt{SHA224\_Digest} (28) or \texttt{SHA256\_Digest} (32).

\lstcap{Example 1: SHA-256 of a message.}
\begin{lstlisting}
with ESP32S3.SHA; use ESP32S3.SHA;
procedure SHA_256_Ex is
   Msg    : constant Byte_Array := (16#61#, 16#62#, 16#63#);  --  "abc"
   Digest : constant SHA256_Digest := Hash_256 (Msg);
begin
   null;   --  Digest is the FIPS-180 vector for "abc"
end SHA_256_Ex;
\end{lstlisting}

\lstcap{Example 2: SHA-1 (e.g. for a legacy checksum).}
\begin{lstlisting}
with ESP32S3.SHA; use ESP32S3.SHA;
procedure SHA_1_Ex is
   Msg : constant Byte_Array := (16#61#, 16#62#, 16#63#);
   D1  : constant SHA1_Digest := Hash_1 (Msg);
begin
   null;
end SHA_1_Ex;
\end{lstlisting}

The accelerator is a shared resource serialised by a protected object, so
concurrent \texttt{Hash} calls from different tasks are safe.

\section{AES: hardware block cipher}\index{AES driver}

\begin{lstlisting}
subtype Key_128 is Key_Bytes (0 .. 15);   subtype Key_256 is Key_Bytes (0 .. 31);
function Encrypt_ECB (Key : Key_Bytes; Plain  : Block) return Block
   with Pre => Supported_Key (Key);        --  16- or 32-byte keys only
function Decrypt_ECB (Key : Key_Bytes; Cipher : Block) return Block
   with Pre => Supported_Key (Key);
\end{lstlisting}

\textbf{Parameters.}
\begin{description}[style=nextline]
\item[\texttt{Key}] the cipher key (\texttt{Key\_Bytes}); its \emph{length} selects
  the cipher: 16 bytes (\texttt{Key\_128}) for AES-128, 32 bytes
  (\texttt{Key\_256}) for AES-256. The \texttt{Supported\_Key} precondition rejects
  any other length (there is no AES-192 on this silicon).
\item[\texttt{Plain} / \texttt{Cipher}] one 16-byte \texttt{Block} to transform; the
  result is a \texttt{Block}.
\end{description}

\lstcap{Example 1: AES-128 encrypt then decrypt one block.}
\begin{lstlisting}
with ESP32S3.AES; use ESP32S3.AES;
procedure AES_128_Ex is
   Key    : constant Key_128 := (16#00#, 16#01#, 16#02#, 16#03#, 16#04#, 16#05#,
                                 16#06#, 16#07#, 16#08#, 16#09#, 16#0A#, 16#0B#,
                                 16#0C#, 16#0D#, 16#0E#, 16#0F#);
   Plain  : constant Block := (others => 16#11#);
   Cipher : constant Block := Encrypt_ECB (Key, Plain);
   Back   : constant Block := Decrypt_ECB (Key, Cipher);   --  Back = Plain
begin
   null;
end AES_128_Ex;
\end{lstlisting}

\lstcap{Example 2: AES-256 (a 32-byte key selects the cipher automatically).}
\begin{lstlisting}
with ESP32S3.AES; use ESP32S3.AES;
procedure AES_256_Ex is
   Key256 : constant Key_256 := (others => 16#2B#);
   Plain  : constant Block := (others => 16#11#);
   C      : constant Block := Encrypt_ECB (Key256, Plain);
begin
   null;
end AES_256_Ex;
\end{lstlisting}

The key parameter is one unconstrained type; its \emph{length} (16 or 32 bytes)
chooses AES-128 or AES-256, and the \texttt{Supported\_Key} precondition rejects
anything else: the contract makes the unsupported AES-192 (absent on this
silicon) a compile/run-time error rather than a silent fallback.
